\addchap{KẾT LUẬN VÀ HƯỚNG PHÁT TRIỂN}\label{chap:conclusion}
\addsec{Kết luận}
Đề tài đã nghiên cứu và phát triển thành công nguyên mẫu hệ thống ứng dụng trí tuệ nhân tạo hỗ trợ tiếp nhận, xử lý và điều phối thông tin cứu hộ bão lũ trong điều kiện mạng gián đoạn. Hệ thống tích hợp toàn diện ba phân hệ: ứng dụng di động Flutter hỗ trợ suy luận AI tại biên, máy chủ FastAPI tích hợp giải thuật gom cụm sự cố và bàn điều phối Web GIS trực quan hóa bản đồ thời gian thực. Các cơ chế lưu trữ cục bộ, gửi dữ liệu thích ứng và khử trùng lặp giao dịch đã chứng minh khả năng đảm bảo tính toàn vẹn thông tin cứu nạn trong môi trường viễn thông khắc nghiệt.

Đối với mô hình thị giác máy tính, kiến trúc MobileNetV3-Large đạt độ chính xác Accuracy \NewImageAccuracy\% và macro-F1 \NewImageMacroFOne\% trên tập ảnh kiểm tra độc lập. Quá trình tối ưu hóa mô hình sang các định dạng thực thi biên (ONNX Runtime, ExecuTorch) duy trì sự đồng thuận dự đoán cao, đồng thời kỹ thuật lượng tử hóa giúp giảm đáng kể kích thước mô hình, đáp ứng yêu cầu triển khai hiệu quả trên thiết bị di động.

Mô hình ước lượng mức độ khẩn cấp dựa trên các đặc trưng nhân khẩu và dấu hiệu lâm sàng cơ bản cung cấp cơ chế tính điểm minh bạch, có thể giải thích được và dễ dàng cấu hình, hỗ trợ đắc lực cho khâu tiền lọc dữ liệu cứu hộ tại hiện trường.

Về mặt giải thuật, các thực nghiệm đối chứng đa dạng trên đồ thị không gian--thời gian đã làm sáng tỏ ưu thế cũng như giới hạn của các phương pháp phân cụm cộng và tích (RQ1), đánh giá tính ổn định của hàm điểm ưu tiên đa tiêu chí trước các biến thể dữ liệu trùng lặp (RQ2), và phân tích định lượng hệ quả của các sai số phân hoạch đối với bài toán lập lịch điều phối nguồn lực cứu trợ (RQ3).

Các đóng góp chính của đề tài bao gồm:
\begin{enumerate}
\item Xây dựng hoàn chỉnh nguyên mẫu hệ sinh thái phần mềm cứu nạn phân tán từ ứng dụng di động biên đến máy chủ điều phối Web GIS;
\item Đề xuất và kiểm chứng giải thuật gom cụm đồ thị không gian--thời gian và hàm ưu tiên cứu nạn đa tiêu chí;
\item Thiết kế phương thức truyền nhận thích ứng và lưu trữ bền vững hỗ trợ mạng gián đoạn, góp phần nâng cao tính chủ động và hiệu quả tác chiến trong công tác cứu nạn thiên tai.
\end{enumerate}

\addsec{Hướng phát triển}
Trên cơ sở các kết quả đạt được, các hướng phát triển tiếp theo của đề tài bao gồm:
\begin{itemize}
\item Hoàn thiện giao tiếp ngữ nghĩa dữ liệu: tích hợp sâu hơn mô hình ngôn ngữ tự nhiên tiếng Việt chuyên biệt cho tin nhắn cứu nạn, nâng cao độ chính xác trích xuất thực thể và dấu hiệu nguy cấp từ văn bản tự do;
\item Mở rộng thử nghiệm phần cứng: triển khai thử nghiệm trên đa dạng các dòng điện thoại Android và iOS thực tế, đánh giá chi tiết mức độ tiêu thụ pin và độ trễ suy luận trong nhiều điều kiện môi trường;
\item Tích hợp mạng truyền thông chịu trễ (DTN): phát triển giao thức chuyển tiếp dữ liệu cơ hội qua kết nối cự ly ngắn (BLE, Wi-Fi Direct) kết hợp các trạm thu phát cơ động hoặc thiết bị bay không người lái (UAV) phục vụ các khu vực bị chia cắt cô lập hoàn toàn;
\item Thử nghiệm thực địa diện rộng: phối hợp cùng các lực lượng cứu hộ chuyên trách và chính quyền địa phương để diễn tập, đánh giá mức độ khả dụng và hoàn thiện quy trình nghiệp vụ cứu nạn trong thực tế bão lũ.
\end{itemize}
